\section{Reduced fibers and descent by a finite flat norm}
\label{g:reduced-flat}

The limiting fiber norm in Lemma~\ref{g:moments} is defined first on
an open subset of the base.  To extend it as a bounded metric, we need a
line whose rational section has the same growth at the boundary.  This
section constructs that line after an alteration with reduced fibers,
then returns it to the original function field by a finite flat norm.
The alteration is used to control a section on every fiber; the final
base is birational to the original base.

We begin with an algebraic construction that explains the role of
reduced fibers.  All schemes in the next two lemmas are of finite type
over \(\mathbb C\).  A modification of an integral base means a proper
birational morphism, and all families over a modified base are pulled
back.

\begin{lemma}[Extending a generic line of sections]
\label{g:flat-section-line}
Let \(g:W\to T\) be a projective flat morphism with geometrically reduced
fibers, where \(T\) is a smooth integral projective variety.  Let \(P\)
be a line bundle on \(W\), and let
\(0\ne\tau\in H^0(W_{\eta},P_{\eta})\), where
\(\eta=\operatorname{Spec}\mathbb C(T)\).
After a smooth projective modification of \(T\), there are a line bundle
\(J\), a rational section \(\eta_J\) of \(J\), and a homomorphism
\[
                 a:g^*J\longrightarrow P
\]
such that \(a(g^*\eta_J)=\tau\) generically and
\(a|_{W_t}\ne0\) for every \(t\in T\).
\end{lemma}

\begin{proof}
Choose a sufficiently high
power \(A\) of a relatively ample line bundle so that
\[
 E_A=g_*A,\qquad F_A=g_*(A\otimes P)
\]
are vector bundles, their formation commutes with arbitrary base change,
and \(g^*E_A\to A\) is surjective.  Relative Serre vanishing and flatness
supply these assertions simultaneously.  Multiplication by \(\tau\)
defines a nonzero element of
\(\operatorname{Hom}(E_A,F_A)_{\eta}\).
Take the closure of the associated rational section of the projective
bundle of lines in \(\operatorname{Hom}(E_A,F_A)\), and resolve this
closure.  The resulting base modification carries a line subbundle
\[
                 J\hookrightarrow\operatorname{Hom}(E_A,F_A).
\]
In particular its map to each fiber of this Hom bundle is nonzero.

The induced map
\(g^*J\otimes g^*E_A\to A\otimes P\) factors through
\(g^*J\otimes A\).  Indeed, the kernel of
\(g^*E_A\to A\) is locally free.  Its composite with this map vanishes
over the generic base point, where the map is multiplication by
\(\tau\).  It therefore vanishes everywhere: flatness over the integral
base makes the generic part schematically dense, and the target is
locally free.  Dividing by \(A\) gives \(a:g^*J\to P\).
If its restriction to some fiber were zero, the corresponding map
\((E_A)_t\to(F_A)_t\) would be zero, contrary to the line subbundle
construction.  The prescribed generic equality defines \(\eta_J\).

\end{proof}

\begin{lemma}[The fiber norm has a positive lower bound]
\label{g:flat-supremum}
In the situation of Lemma~\ref{g:flat-section-line}, give \(J\) and
\(P\) continuous positive Hermitian metrics.  For
\(a:g^*J\to P\), put
\[
                  c(t)=\sup_{w\in W_t}|a(w)|^2.
\]
There are constants \(0<c_-\le c_+<\infty\) with
\(c_-\le c(t)\le c_+\) for all \(t\in T\).
\end{lemma}

\begin{proof}
Properness gives the upper bound.  On a reduced fiber, a nonzero section
of a line bundle is nonzero at a point in its smooth locus.  Choose such
a point over \(t\).  Flatness and smoothness of the fiber at that point
make \(g\) smooth there.  After shrinking, a local analytic section of
\(g\), together with continuity of the metric, supplies a positive lower
bound on a neighborhood of \(t\).  Finitely many such neighborhoods
cover the compact base.  Reducedness is used precisely here: a nonzero
section supported only in the nilpotent structure of a fiber would not
supply a point at which its norm is positive.
\end{proof}

\begin{remark}[Proper flat families]
The conclusion of Lemma~\ref{g:flat-section-line} also holds when
\(g\) is proper and flat without being projective. To prove this
variant, replace the Hom bundle construction by the complex computing
proper flat cohomology. The complex \(Rg_*P\) is perfect and commutes
with derived base change \cite[Tag~0DJT]{StacksProject}. On the smooth
projective base it has a global bounded vector bundle representative,
by the resolution property \cite[Tags~0F87 and 0F8F]{StacksProject}.
It can be chosen in degrees \(0,\ldots,N\): the derived fiber
cohomology has no negative degrees, so the negative part of a bounded
representative is successively removed by locally split injections and
their vector bundle cokernels. Write the resulting complex as
\[
              K^0\xrightarrow{d^0}K^1\longrightarrow\cdots.
\]
The vector \(\tau\) is a nonzero element of
\(\ker(d^0_{\eta})\subset K^0_{\eta}\). Modify the base to extend its
line to a subbundle \(J\subset K^0\). The composite \(J\to K^1\)
vanishes generically and hence everywhere. Thus
\(J\to\ker d^0=g_*P\) gives the required evaluation map.
On each fiber, the subbundle gives a nonzero vector in
\(\ker d^0_t=H^0(W_t,P_t)\); there is no incoming differential in
negative degree. This proves nonvanishing on every fiber. The
projective construction suffices for the transfer below.
\end{remark}

\subsection{The norm line on the original field}

We now use these lemmas to construct the two base metrics and the map
that will return sections to the source.  The anticanonical metric and
the bounded metric in the statement are separate data.

\begin{proposition}[Reduced-flat transfer]
\label{g:reduced-flat-transfer}
Let \(X,V,Y\) be smooth integral projective varieties, let
\(\pi:V\to X\) be birational, and let \(f:V\to Y\) be surjective with
connected fibers.  Let \(D\) be a reduced effective divisor on \(X\),
let \(c\ge0\) be an integer, and let \(L\) be a line bundle.
Suppose that \(L\) has a semipositive metric \(h_L\) with locally bounded
weights and that \(M=-K_X+cD\) has a semipositive singular metric
\(h_M\).  In local coordinates and a compatible frame of \(M\), write
\[
              d\mu=|d|^{2c}e^{-\varphi_M}\,dV_X,
\]
where \(d\) is an equation of \(D\), and assume \(\mu(X)<\infty\).
Put \(E=K_{V/X}\).  Assume that for some integer \(a>0\) there is a
nonzero rational section \(\tau\) of \(a\pi^*L\), regular over a dense
base open subset, and that
\begin{equation}
 \operatorname{rk} f_*\mathcal O_V\bigl(E+c\pi^*D+pa\pi^*L\bigr)=1
                \qquad\text{for every integer }p\ge0.
 \label{g:flat-adjoint-ranks}
\end{equation}
If \(\dim Y>0\), there are a smooth projective modification
\(Y'\to Y\), a line bundle \(B\) on \(Y'\), a reduced effective divisor
\(D'\), and integers \(d_0>0\), \(c'\ge0\) with the following properties.
The line \(B\) has a semipositive metric with locally bounded weights.
The line \(-K_{Y'}+c'D'\) has a semipositive metric inducing a finite
measure.  On a smooth resolution
\(\pi':V'\to X\), \(f':V'\to Y'\) of the main transform there is a
nonzero regular homomorphism
\begin{equation}
                  (f')^*B\longrightarrow d_0a(\pi')^*L,
 \label{g:flat-output-map}
\end{equation}
and
\begin{equation}
       \operatorname{Supp}\bigl(\pi'_* (f')^*D'\bigr)
                            \subseteq\operatorname{Supp}D.
 \label{g:flat-pole-support}
\end{equation}
Consequently, for every \(m>0\), a nonzero rational section of \(mB\)
with poles only on \(D'\) gives a nonzero rational section of
\(md_0aL\) with poles only on \(D\).

There is also a smooth invariant version.  In that version \(D=0\),
both metrics are smooth, an algebraic torus acts compatibly on \(X,V\)
and the line bundles, its action on \(Y\) is trivial, and \(\tau\) is
invariant.  A Zariski dense compact subgroup preserves the metrics.
Replace the ranks in \eqref{g:flat-adjoint-ranks} by the ranks of the
invariant parts, still for every \(p\ge0\).  Then the transferred section
on \(X\) is invariant and regular.  Anticanonical bundles in this version
carry their natural linearizations.
\end{proposition}

\begin{proof}
On a dense smooth base open, use coordinates \(y\) and write \(I_p(y)\)
for the density of
\[
              f_*\bigl(|\tau|_{(\pi^*h_L)^a}^{2p}\,\pi^*\mu\bigr).
\]
The pullback measure includes the Jacobian of \(\pi\).  The adjoint
section is \(s_Es_{c\pi^*D}\tau^p\); its fiberwise squared norm is
\(I_p\), after trivializing the canonical base factor.
Lemma~\ref{g:moments} applies to
\eqref{g:flat-adjoint-ranks}, including \(p=0\).  It gives psh
representatives of \(-\log I_0\) and
\begin{equation}
 u=-\log S,\qquad
 S(y)=\operatorname*{ess\,sup}_{V_y}|\tau|^2_{(\pi^*h_L)^a}.
 \label{g:flat-limit-norm}
\end{equation}
The essential supremum uses the fiber density of \(\pi^*\mu\).
For almost every smooth fiber that density is positive almost everywhere.
Since \(h_L\) has bounded weights, \(S\) is comparable to the ordinary
supremum computed with a fixed smooth reference metric.  The comparison
constants are uniform on the compact total space wherever the rational
section is defined. Fix this smooth reference metric on \(L\), and use
its pullbacks on every model below. In the invariant case the same lemma uses
the holomorphic orthogonal projection to the invariant direct summand
for the smooth coefficient metric.  Positivity of that quotient metric
is the additional input required in this case.

\emph{Constructing a line on an alteration.}
Apply weak semistable reduction of Abramovich--Karu to \(f\)
\cite[Theorem~0.3 and Definition~0.1, arXiv version~1]{AK2000}.  After a projective generically finite
alteration \(T\to Y\) and modifications, it gives a projective family
\(g:W\to T\) with geometrically reduced fibers, with \(T\) smooth
projective, and a morphism \(W\to V\) over \(Y\).  Its equidimensional
toroidal source is Cohen--Macaulay, so it is flat over the smooth base.
The source \(W\) need not be smooth.  We keep this flat family when
making further base modifications; resolving its source would not
preserve the flatness needed below.  The generic fiber of \(f\) is
geometrically integral, because it is smooth and connected in
characteristic zero, as required for this application of reduction.
The two families fit into the diagram
\[
 \begin{array}{ccc}
 W&\longrightarrow&V\\
 {\scriptstyle g}\downarrow&&\downarrow{\scriptstyle f}\\
 T&\longrightarrow&Y.
 \end{array}
\]

Let \(P\) be the pullback to \(W\) of \(a\pi^*L\).
Lemma~\ref{g:flat-section-line} gives \(J\), \(\eta_J\), and
\(g^*J\to P\) on a further modified base.  The family remains flat with
reduced fibers.  On the generic base open the map \(W\to V\) is
surjective on the fibers under consideration. Give \(J\) a continuous
positive reference metric and let \(c(t)\) be the supremum norm of the
map \(g^*J\to P\), measured in these reference metrics. Since this map
sends \(g^*\eta_J\) to \(\tau\), we have
\[
 \sup_{W_t}|\tau|_{\mathrm{ref}}^2
     =|\eta_J(t)|_{\mathrm{ref}}^2c(t),
 \qquad c_-\leq c(t)\leq c_+
\]
by Lemma~\ref{g:flat-supremum}. Surjectivity onto the original fiber
identifies the left side with \(\sup_{V_y}|\tau|_{\mathrm{ref}}^2\).
Comparison with the bounded metric \(h_L\) therefore gives
\begin{equation}
                  |\eta_J(t)|^2\asymp S(y),\qquad t\mapsto y,
 \label{g:flat-upstairs-comparison}
\end{equation}
for almost every parameter in that open.  Here and below comparison
means upper and lower positive constant bounds in local reference
metrics, with constants that remain bounded toward the omitted locus.
The fixed reference metrics are pulled back under all subsequent base
changes, and the same bounds on \(c(t)\) continue to hold.
No rational connectedness of \(T\) is asserted or needed.

\emph{Returning to the original function field.}
By flattening by modification, choose a smooth projective modification
\(Y'\to Y\) on which both the main transform \(V_0\to Y'\) of \(V\)
and the main transform \(q:T'\to Y'\) of \(T\) are flat
\cite{RaynaudGruson1971}.  The second map is finite: it is projective,
flat and has zero-dimensional fibers.  Its degree
\[
                  d_0=\operatorname{rk}q_*\mathcal O_{T'}
\]
is positive and constant.  The first map is equidimensional of the
original relative dimension.  Subsequent smooth modifications of the
base preserve these properties by flat base change.  The main
transforms map to \(T\) and \(V\), respectively, so \(J\) and its
rational section pull back to \(T'\).  The base \(Y'\) still has the
function field \(\mathbb C(Y)\).
The resulting diagram is
\[
 \begin{array}{ccccc}
 T'&\xrightarrow{\ q\ }&Y'&\longleftarrow&V_0\\
 \downarrow&&\downarrow&&\downarrow\\
 T&\longrightarrow&Y&\longleftarrow&V.
 \end{array}
\]
Here both maps to \(Y'\) are flat, and \(q\) is finite. The line
\(J\) lives on \(T'\); its norm will live on \(Y'\).

Define the actual norm line and its rational section by
\begin{equation}
 B=\operatorname{Nm}_q(J)
   :=\det(q_*J)\otimes\det(q_*\mathcal O_{T'})^{-1},
 \qquad \beta=\operatorname{Nm}_q(\eta_J).
 \label{g:flat-norm-definition}
\end{equation}
Both direct images in this formula are locally free of rank \(d_0\).
To describe the section, view multiplication by \(\eta_J\) as a rational
map \(q_*\mathcal O_{T'}\to q_*J\) and take its determinant.
Locally a line bundle can be trivialized near a finite inverse fiber;
changing the frame by a unit changes its norm frame by the determinant
of multiplication by that unit.  This proves the compatibility of
\eqref{g:flat-norm-definition} with changes of frame and arbitrary base
change.  The construction uses finite local freeness, and does not
require \(T'\) to be smooth or \(q\) to be unramified
\cite[Tags~0BUT, 0BCX and 0BD2]{StacksProject}.

Over the generically unramified locus the norm is the product over the
\(d_0\) sheets.  Applying \eqref{g:flat-upstairs-comparison} on those
sheets gives
\begin{equation}
                         |\beta(y)|^2\asymp S(y)^{d_0}.
 \label{g:flat-norm-comparison}
\end{equation}
These constants remain bounded at the branch locus.  Indeed, choose a
frame for \(J\) near the finite inverse fiber.  Its continuous reference
norm is bounded above and away from zero there, and its norm is a
nonvanishing frame of \(B\).  Products of the finitely many reference
norms are therefore comparable to any fixed reference norm on \(B\).
This also explains why the exponent is multiplied by \(d_0\).

Define a metric on \(B\), initially where \(\beta\) is a nonzero regular
frame, by \(|\beta|^2=S^{d_0}\).  If \(\beta=b e\) in a local
holomorphic frame \(e\), its weight is
\[
                    \psi=d_0u+\log|b|^2.
\]
It is psh on that open.  Equation~\eqref{g:flat-norm-comparison} bounds
\(\psi\) above and below locally even toward the omitted locus.  The
bounds hold for the psh representatives, by subaveraging and upper
semicontinuity.  Removable singularities across analytic sets therefore
extend \(\psi\) to a locally bounded psh weight on all of \(Y'\).

Resolve \(V_0\) to \(V'\), with maps \(\pi'\) and \(f'\) as in the
statement.  The rule
\begin{equation}
                          \beta\longmapsto\tau^{d_0}
 \label{g:flat-section-rule}
\end{equation}
defines a rational homomorphism in \eqref{g:flat-output-map} over the
original field. It extends regularly. Indeed, write
\(\beta=b e_B\) and \(\tau=t e_L^{\otimes a}\) in local holomorphic
frames, with the fixed reference norms bounded above and below. The
coefficient of this map is \(t^{d_0}/(b\circ f')\). On almost every
fiber in the good open, \eqref{g:flat-norm-comparison} and the definition
of the essential supremum give
\[
 \left|\frac{t^{d_0}}{b\circ f'}\right|^2
 \leq C\frac{S(f'(\,\cdot\,))^{d_0}}
                  {|\beta\circ f'|_{\mathrm{ref}}^2}
 \leq C'.
\]
The bound first holds almost everywhere on these fibers and then
everywhere on their good locus by continuity of the rational coefficient.
A rational coefficient with a pole along a prime divisor would be
unbounded on open subsets approaching its general point; those subsets
meet these fibers.  This excludes every divisorial pole, and normality
extends the map across codimension two.  In the invariant case the rule
\eqref{g:flat-section-rule} is invariant as a rational identity on \(X\).
The auxiliary alterations and final resolution need not be equivariant.

We have constructed the bounded base metric and its map back to \(X\).
On the new base, \(S\) denotes the same generic fiber supremum, whereas
\(I_0\) now denotes the density of \((f')_*(\pi')^*\mu\) in the chosen
coordinates on \(Y'\).  The latter transforms as a volume density,
including the base Jacobian.  Off the exceptional locus of the base
modification, this changes \(-\log I_0\) by a pluriharmonic coordinate
term and retains its plurisubharmonicity.
It remains to extend the volume metric and ensure that all its correction
primes disappear, or land in \(D\), upon return to \(X\).

\emph{The correction divisor and every component of its pullback.}
Among the boundary prime divisors of the good open in \(Y'\), let \(D'\)
be the reduced sum of those primes \(Q\) for which every prime of
\(V'\) dominating \(Q\) is either \(\pi'\)-exceptional or maps into
\(D\).  This definition first addresses components that dominate base
primes.  It gives \eqref{g:flat-pole-support} for \emph{all} components
of the total pullback because of the equidimensional model \(V_0\).
To see this, put \(n=\dim X\) and \(r=n-\dim Y'\).  If a prime
\(R\subset V'\) maps into a set \(A\subset Y'\) of codimension at
least two, then
\[
       \dim(V_0\times_{Y'}A)\le\dim A+r\le n-2.
\]
The image of \(R\) in \(V_0\), hence in \(X\), has dimension at most
\(n-2\), since \(V_0\to X\) is a morphism.  Thus \(R\) is
\(\pi'\)-exceptional.  Every other component of \((f')^*D'\) dominates
a component of \(D'\) and is covered by its definition.  This proves
\eqref{g:flat-pole-support}, including the divisors introduced by the
last resolution.

At a general point of a boundary prime \(Q\), choose a component
\(R\subset(f')^*Q\) dominating it, and write
\[
 k=\operatorname{ord}_R(f')^*Q>0,\qquad
 b=\operatorname{ord}_R\bigl(K_{V'/X}+c(\pi')^*D\bigr)\ge0.
\]
The psh weight of \(h_M\) is locally bounded above.  Hence the pullback
measure is bounded below by a positive constant times
\(|z_1|^{2b}\) times coordinate volume near a general point of \(R\).
Coordinates can be chosen there with \(t_1=z_1^k\), where \(Q=(t_1=0)\),
and with the other base coordinates among the remaining source
coordinates.  Integration in a fixed small fiber polydisc and change of
variables give
\begin{equation}
                    I_0(t)\ge C|t_1|^{2((b+1)/k-1)}
 \label{g:flat-volume-lower}
\end{equation}
almost everywhere near a general point of \(Q\).  If \(Q\not\subset D'\),
choose \(R\) nonexceptional over \(X\) and not mapping into \(D\).
Then \(b=0\); the exponent in \eqref{g:flat-volume-lower} is nonpositive,
so in particular \(I_0\) has a positive constant lower bound there.

There are only finitely many boundary primes.  Choose an integer
\(c'\ge0\) at least as large as each positive exponent
\((b+1)/k-1\) needed for a component of \(D'\).  If \(d'\) is an equation
of \(D'\), the psh function
\[
                     -\log I_0+c'\log|d'|^2
\]
is locally bounded above at the generic points of all boundary primes.
It extends across them, and then across the remaining analytic subset
of codimension at least two by the Hartogs extension property for psh
functions.  These are the weights of a semipositive metric on
\(-K_{Y'}+c'D'\).  Dividing its metric by the divisor metric of \(c'D'\)
recovers the measure with density \(I_0\).  It has finite mass because it
is the pushforward of the original finite measure.

Finally take a nonzero rational section of \(mB\) with poles on \(D'\),
pull it back, and apply the \(m\)-th power of
\eqref{g:flat-output-map}.  The only possible poles lie in
\((f')^*D'\).  Equation~\eqref{g:flat-pole-support} shows that its
pushdown has poles only on \(D\), and it is nonzero at the generic point.
In the smooth invariant case this rational section is invariant by
\eqref{g:flat-section-rule}; since \(D=0\), normality makes it a global
regular section on \(X\).
\end{proof}

\begin{remark}[Valuations of the norm]
\label{g:flat-norm-valuations}
The finite norm itself creates no polar prime away from the images of
polar primes upstairs.  At a prime \(Q\subset Y'\), pass to the discrete
valuation ring \(\mathcal O_{Y',Q}\) and normalize its finite generic
extension.  If \(v_i\) are the extension valuations and \(f_i\) their
residue degrees, then for a nonzero rational function \(z\),
\[
                \operatorname{ord}_Q\operatorname{Nm}(z)
                             =\sum_i f_i v_i(z).
\]
This is the determinant of multiplication by \(z\), computed by lengths
of lattices after clearing denominators.  Applied in local frames, it
gives the same formula for a normed rational section.  For the section
returned to \(X\), the stronger statement needed is
\eqref{g:flat-pole-support}: it also treats components lying over
codimension-two sets in the base.
\end{remark}

\begin{remark}[What is descended]
The determinant construction in \eqref{g:flat-norm-definition} descends
a line on a possibly ramified, nonsmooth finite flat cover and changes
\(a\) to \(d_0a\).  The finite-group quotient construction in
\cite[Subsection ``Descending the alteration and its section'']{CCohomologicalTransfer}
instead first kills stabilizer actions by
a tensor power and then uses a function-field norm on a normal target.
These are different model constructions.  They are also distinct from
the finite \emph{\'etale} section norm of a line already pulled back from
the original variety.  None of the three permits induction on the
altered base \(T\) here; the induction below takes place on \(Y'\),
whose function field remains \(\mathbb C(Y)\).
\end{remark}

\subsection{An approximation proof of the slope input}

For the rationally connected application of the reduced-flat
construction, the required slope inequality also has an approximation
proof.  Its distinctive step is to convert finite weighted volume into
a klt boundary on a resolution.  We give that step, rather than infer
it from the existence of the descended metric.

\begin{proposition}[Floor approximation for foliation determinants]
\label{g:reduced-flat-slope}
Let \(X\) be smooth projective, let \(D\) be a reduced effective divisor,
and let \(c\ge0\) be an integer. Suppose that \(M=-K_X+cD\) has a
semipositive singular metric with finite induced measure
\(|d|^{2c}e^{-\varphi}\,dV_X\). Every algebraically integrable saturated
foliation \(\mathcal F\subset T_X\) of positive rank and corank satisfies
\begin{equation}
                         K_{\mathcal F}+M
                          \quad\hbox{is pseudoeffective}.
 \label{g:flat-foliation-psef}
\end{equation}
\end{proposition}

\begin{proof}
Its leaves define a rational fibration.  Fix a sufficiently ample line
bundle \(H\), independent of the integer \(l\ge2\).  Nadel vanishing
and Castelnuovo--Mumford regularity make
\[
               \mathcal O_X(lM+H)\otimes\mathcal J(l\varphi)
\]
globally generated.  The local Bergman approximation inequality gives,
for its local base ideal generators \(g_j\),
\begin{equation}
                          e^{l\varphi}\le C_l\sum_j|g_j|^2.
 \label{g:flat-approximation}
\end{equation}
The fixed ample twist can be chosen by applying Nadel vanishing after
subtracting \(1,\ldots,\dim X\) times a fixed very ample line bundle.
The local inequality follows from point extension in the weighted
holomorphic \(L^2\) space; generators of the multiplier ideal represent
the singularities of that space.  These are the usual Bergman
approximation and global generation constructions of Demailly
\cite[Theorems~0.3 and 1.2]{Demailly2015}.

Let \(\rho:U\to X\) resolve the rational fibration and principalize this
base ideal, with its divisor \(R_l\), and let \(E_U=K_{U/X}\).  Take a
simultaneous log resolution with \(\rho^*D\).  Finite mass and
\eqref{g:flat-approximation} imply, along every prime on \(U\),
\begin{equation}
                         \lfloor R_l/l\rfloor
                                  \le E_U+c\rho^*D.
 \label{g:flat-floor-bound}
\end{equation}
Indeed the comparison makes
\(|s_{E_U+c\rho^*D}|^2|s_{R_l}|^{-2/l}\) locally integrable.
At a prime with respective orders \(b\) and \(r\), integrability says
\(b-r/l>-1\); since \(b\) is integral this is exactly
\(\lfloor r/l\rfloor\le b\).

Pull back the linear subsystem
\[
 H^0\bigl(X,\mathcal O_X(lM+H)\otimes\mathcal J(l\varphi)\bigr)
                  \subseteq H^0(X,lM+H)
\]
to \(U\), and remove its fixed divisor \(R_l\).  Let \(G_l\) be a
general member of the resulting basepoint-free system, and put
\[
           \Delta_l=R_l/l-\lfloor R_l/l\rfloor+G_l/l.
\]
Bertini and the log resolution make \(\Delta_l\) an effective klt
boundary, also klt on the general fiber of the resolved fibration.
Moreover,
\begin{equation}
 K_U+\Delta_l\sim_{\mathbb Q}
       E_U+c\rho^*D-\lfloor R_l/l\rfloor+\rho^*H/l.
 \label{g:flat-adjoint-approximation}
\end{equation}
By \eqref{g:flat-floor-bound}, a positive divisible multiple of the
right side restricts to an effective divisor on the general fiber.
One may take a general effective representative of \(H\), which does
not contain that fiber.

The relative pseudoeffectivity theorem for an algebraically integrable
foliation now gives
\(K_{\mathcal F_U}+\Delta_l\) pseudoeffective on \(U\), where
\(\mathcal F_U\) is the induced foliation
\cite[Proposition~3.1, arXiv version~4]{Ou2023}; its relative positivity
input is due to Druel \cite[Proposition~4.1]{Druel2017}.
The hypotheses used are the effective boundary, log canonicity of its
general-fiber restriction, and the general-fiber nonvanishing in
\eqref{g:flat-adjoint-approximation}.  Push down to \(X\), and add back
the effective divisor \(\rho_*\lfloor R_l/l\rfloor\).  Since foliation
determinants agree away from the exceptional locus, this proves that
\(K_{\mathcal F}+M+H/l\) is pseudoeffective.  Its limit as
\(l\to\infty\) proves \eqref{g:flat-foliation-psef}.

\end{proof}

Apply Lemma~\ref{g:bt:slope-principle} with \(P=cD\).
For every nonzero movable class \(\alpha\) on a positive-dimensional
\(X\) with \(D\cdot\alpha=0\), the preceding proposition gives
\(\mu_{\alpha,\min}(T_X)\ge0\). Movable tensor semistability then gives
nonpositive maximal slope for every cotangent tensor power and its
exterior powers. Thus every coherent subsheaf of
\((\Omega_X^1)^{\otimes k}\), or of \(\Omega_X^p\), has nonpositive
\(\alpha\)-degree. For degree zero this follows directly from the
inclusion in \(\mathcal O_X\). The floor construction above is an
alternative proof of the foliation input to the general weighted
estimate, Proposition~\ref{g:bt:weighted-slope}.

\subsection{An alternative proof of allowed-pole nonvanishing}
\label{g:rc-allowed-poles}

The reduced-flat transfer preserves the original base field.  This is
what permits rational connectedness to pass to the base on which we
induct.  The following proof records how the common field and
determinant constructions meet every rank hypothesis of
Proposition~\ref{g:reduced-flat-transfer}.

We prove the rationally connected case of Theorem~\ref{g:allowed}
using reduced-flat transfer. First reduce its arbitrary effective
integral boundary \(D_0\) to a multiple of a reduced divisor. Put
\(D=\operatorname{Supp}(D_0)_{\mathrm{red}}\), choose an integer
\(c\geq0\) with \(cD-D_0\geq0\), and tensor the metric on
\(-K_X+D_0\) with the canonical divisor metric of \(cD-D_0\).
The new metric on \(-K_X+cD\) is semipositive, and its distinguished
measure is unchanged: the squared divisor factor cancels the added
metric factor. A rational section of \(mL\), \(m>0\), with poles only
on \(D\) belongs to \(H^0(X,mL+rD_0)\) for some sufficiently large
\(r\geq0\).
The zero boundary is included by taking \(D=0\) and \(c=0\).

It therefore suffices to take \(X\) smooth projective and rationally
connected, \(D\) reduced effective, and \(L\) a line bundle that
has a semipositive metric with locally bounded weights. Suppose
\(-K_X+cD\) has a semipositive metric with finite induced measure for
an integer \(c\ge0\). We seek a nonzero section of \(mL+rD\) with
integers \(m>0\) and \(r\ge0\). A point is included.

\begin{proof}
Induct on \(n=\dim X\); the assertion for a point is immediate.
Suppose the conclusion fails.  The class \(L+D\) is pseudoeffective
and not big, so cone duality gives a nonzero movable class \(\alpha\)
with
\[
                       L\cdot\alpha=D\cdot\alpha=0
\]
\cite{BDPP2013}.  Proposition~\ref{g:reduced-flat-slope} and the following slope argument apply.
Rational connectedness gives \(H^i(X,\mathcal O_X)=0\) for \(i>0\), so
\(\chi(X,K_X)=(-1)^n\ne0\).  Riemann--Roch makes
\(\chi(X,K_X+mL)\) a nonzero polynomial.  Thus one cohomological degree
occurs with nonzero cohomology for an unbounded sequence of positive
\(m\).  Hard Lefschetz with multiplier ideals supplies twisted
holomorphic forms for that sequence; the multiplier ideals are trivial
because the metric weights on \(mL\) are bounded
\cite{DPS2001}.

Apply the determinant construction of Lemma~\ref{g:determinant}
to these forms.  Its fixed determinant maps nontrivially to a
cotangent tensor power, so its \(\alpha\)-degree is nonpositive by the slope consequence of
Proposition~\ref{g:reduced-flat-slope}.  We obtain a fixed line bundle
\(B_0\) and effective divisors
\begin{equation}
                 N_j\sim jL+B_0,\qquad \alpha\cdot N_j=0
 \label{g:flat-determinant-sequence}
\end{equation}
for an unbounded set of positive integers \(j\).

Let \(K_\alpha\subset\mathbb C(X)\) be the field of functions whose
polar divisors have \(\alpha\)-degree zero.  By
Lemma~\ref{g:cf:null-field}, it is relatively algebraically closed,
its transcendence degree is less than \(n\), and it is represented by
a rational fibration with connected generic fiber.  Resolve it to
\(\pi:V\to X\), \(f:V\to Y\), with \(Y\) smooth projective.
For a general sufficiently ample divisor \(H_Y\) on \(Y\),
\begin{equation}
                         \alpha\cdot\pi_*f^*H_Y=0.
 \label{g:flat-base-degree}
\end{equation}
Indeed a general ratio from \(|H_Y|\) belongs to \(K_\alpha\); the two
general pulled-back divisors have no common component, so its pole
divisor represents the indicated class.

The affine interpolation of Lemma~\ref{g:cf:horizontal-slope} applied
to \eqref{g:flat-determinant-sequence} produces an integer \(a>0\) and
a nonzero rational section \(\tau\) of \(a\pi^*L\) that is regular over
a dense base open.  Explicitly, the linear equivalence
\[
                (k-i)N_j\sim(k-j)N_i+(j-i)N_k
                    \qquad(i<j<k)
\]
has its defining function in \(K_\alpha\).  Consequently the
coefficients of \(\pi^*N_j\) at every horizontal prime are affine in
\(j\); nonnegativity at unbounded \(j\) makes each slope nonnegative.
A difference of two such divisors has no negative horizontal part and
gives \(\tau\).  If \(Y\) is a point, this already gives a global
section of a positive multiple of \(L\).

Suppose therefore that \(\dim Y>0\), and put \(E=K_{V/X}\).
For every integer \(p\ge0\),
\(s_Es_{c\pi^*D}\tau^p\) is a nonzero generic-fiber section of
\(E+c\pi^*D+pa\pi^*L\).  Any finite set of other generic-fiber
sections becomes global after a common sufficiently ample base twist.
The pushdowns of their effective divisors have \(\alpha\)-degree zero,
by \eqref{g:flat-base-degree}, \(\pi_*E=0\), and
\(\alpha L=\alpha D=0\).  Ratios of two such sections therefore
belong to \(K_\alpha\), proving
\eqref{g:flat-adjoint-ranks} for every \(p\), including zero.
This argument establishes generic ranks and requires no simultaneous
base-change assertion for infinitely many sheaves.

Proposition~\ref{g:reduced-flat-transfer} produces \(Y',B,D'\).
The variety \(Y'\) is rationally dominated by \(X\), so it is
rationally connected, and its dimension is smaller than \(n\).
Induction gives a section of \(mB+rD'\) for some \(m>0\), \(r\ge0\).
Transfer returns a nonzero rational section of \(md_0aL\) with poles
only on \(D\), the required contradiction.
\end{proof}

\subsection{An alternative proof for the invariant quotient}
\label{g:flat-invariant-application}

We now prove the rationally connected case of Theorem~\ref{rb:invariant}
by the same reduced-flat method. Let \(F\) be smooth projective and
rationally connected, with \(L=-K_F\) smoothly semipositive, and let an
algebraic torus \(T\) act with the natural anticanonical linearization.
Average a smooth semipositive metric over a compact real form of \(T\).
We seek \(H^0(F,mL)^T\ne0\) for some integer \(m>0\).

\begin{proof}
For the trivial action this is the argument of Section~\ref{g:rc-allowed-poles} with
\(D=0\).  Otherwise take a rational quotient and an equivariant
resolution \(\pi:V\to F\), \(f:V\to Y\).  Its invariant field is
\(\mathbb C(Y)\), and the generic fiber is an orbit closure.  In
particular it is connected.  On the smooth quotient locus, choose
\(\dim V-\dim Y\) generically independent vector fields from the torus
action and wedge them.  Commutativity makes their wedge an invariant
section of \(-K_{V/Y}\).  Tensor it with a rational anticanonical frame
from \(Y\), then multiply by the Jacobian section \(s_E\), where
\(E=K_{V/F}\).  This produces an invariant rational section
\(\tau\) of \(\pi^*L\), regular over a dense base open.

For every \(p\ge0\), the invariant generic-fiber adjoint space for
\(E+p\pi^*L\) contains \(s_E\tau^p\).  Two nonzero invariant sections
have an invariant ratio, hence a ratio in \(\mathbb C(Y)\).
Its invariant rank is exactly one, also at \(p=0\).
Apply the smooth invariant part of
Proposition~\ref{g:reduced-flat-transfer}, using the given smooth metric
for both \(L\) and \(-K_F\).  If \(Y\) is a point, \(\tau\) is already
a global invariant section.  Otherwise the resulting original-field
base \(Y'\) is rationally connected; apply
the argument of Section~\ref{g:rc-allowed-poles} there and transfer its section.
The pole-support statement and \(D=0\) make the resulting invariant
section regular on \(F\).
\end{proof}
